\documentclass[11pt]{article}
\usepackage{amsmath, amssymb}
\usepackage{geometry}
\geometry{margin=1in}

\title{Near-Frontier with Reset (NFR) Planner: Formal Overview}
\author{ARC-AGI Abstraction Navigator Notes}
\date{\today}

\begin{document}
\maketitle

\section{Problem Setting}

Let $\mathcal{G} = (S, A, T)$ denote the agent interaction graph. States $s \in S$ encode masked ARC frames, actions $a \in A$ are constrained to the four arrows, and $T : S \times A \to S \cup \{\bot\}$ is the partial transition function estimated from experience ($\bot$ indicates an unknown successor).

Each play starts at a known level start state $s_0 \in S$. The planner is given:
\begin{itemize}
    \item $S_{\mathrm{disc}} \subseteq S$: discovered states from prior episodes.
    \item $T_{\mathrm{known}} : S_{\mathrm{disc}} \times A \to S_{\mathrm{disc}} \cup \{\bot\}$: the partial transition map recovered from experience (returning $\bot$ when a transition has not yet been seen).
    \item $A(s) \subseteq A$: actions the environment currently allows at state $s$.
    \item $\mathcal{F} = \{s \in S_{\mathrm{disc}} \mid \exists a \in A(s) : T_{\mathrm{known}}(s,a) = \bot\}$, the \emph{frontier} states whose outgoing edges are incompletely observed.
\end{itemize}

During an episode the agent sits in state $s_t$ with available actions $A_t = A(s_t)$. The planner must output $a_t \in A_t \cup \{\mathrm{RESET}\}$ to minimise the distance to a frontier while optionally using reset jumps to $s_0$.

\section{Algorithm}

\subsection{Definition}

Given the current state $s_t$, the available actions $A_t = A(s_t)$, and the level start state $s_0$, the planner executes:
\begin{enumerate}
    \item Run breadth-first search on the known subgraph $\mathcal{G}_t = (S_{\mathrm{disc}}, A, T_{\mathrm{known}})$ from both $s_t$ and $s_0$:
    \begin{align}
    \label{eq:bfs-current}
    D_t^{c}, P_t^{c} &= \mathrm{BFS}(\mathcal{G}_t, s_t),\\
    \label{eq:bfs-reset}
    D_t^{0}, P_t^{0} &= \mathrm{BFS}(\mathcal{G}_t, s_0).
    \end{align}
    Here $D_t^{c}$ and $D_t^{0}$ record shortest-path distances, while the predecessor maps $P_t^{c}$ and $P_t^{0}$ associate each reachable vertex $x$ with the pair $(\mathrm{parent}(x), \mathrm{action}(x))$ describing the final step on the corresponding shortest path.
    \item Compute the reset-augmented cost for every frontier $f \in \mathcal{F}$:
    \begin{equation}
    \label{eq:cost}
    J(f) = \min\{ D_t^{c}(f), 1 + D_t^{0}(f) \}.
    \end{equation}
    \item Select the frontier
    \[
    f^\star = \arg\min_{f \in \mathcal{F}} \big(J(f), D_t^{0}(f), f\big),
    \]
    using lexicographic order on the triple $(J(f), D_t^{0}(f), f)$.
    \item If $J(f^\star)$ is finite and $A_t$ is non-empty, return
    \[
    a_t = \text{first edge on the path recovered from }P_t^{c}\text{ between }s_t\text{ and }f^\star
    \]
    when $D_t^{c}(f^\star) \leq 1 + D_t^{0}(f^\star)$, and return $\mathrm{RESET}$ otherwise.
\end{enumerate}

When $J(f^\star)$ is infinite (no reachable frontier) or $A_t$ is empty, the procedure returns no action.

\subsection{Implementation Notes}

The code performs the following additional checks:
\begin{itemize}
    \item If the frontier loop in step 2 yields no finite-cost candidate, the planner abstains and the caller falls back to its default policy.
    \item If the chosen frontier satisfies $s_t = f^\star$ or the predecessor map does not contain a route, the planner probes the first $a \in A_t$ with $T_{\mathrm{known}}(s_t, a) = \bot$; it abstains when every available action is already known.
    \item The $\mathrm{RESET}$ branch is triggered only when the inequality $D_t^{c}(f^\star) > 1 + D_t^{0}(f^\star)$ holds, i.e.\ restarting from $s_0$ plus one action is strictly cheaper than moving directly from $s_t$.
\end{itemize}

\section{Desirable Properties}

\paragraph{Completeness on the Known Subgraph.} As long as $\mathcal{G}_t$ is connected between $s_t$ and all frontiers, $\mathrm{BFS}$ ensures $D_t^{c}(f)$ equals the true shortest distance in $\mathcal{G}_t$. Therefore the navigation step reaches $f^\star$ in minimal steps without reset.

\paragraph{Reset Optimality.} The cost \eqref{eq:cost} chooses between continuing from $s_t$ or incurring one reset plus the path from $s_0$. This is optimal under the assumption that resetting teleports to $s_0$ in one step and costs a unit time penalty.

\paragraph{Frontier Prioritisation.} Any frontier with strictly smaller $J(f)$ than all others will eventually be chosen. If exploration uncovers new edges, $J$ automatically updates without retuning.

\paragraph{Idempotence of Known Actions.} The planner never replays an action whose successor is already in $T_{\mathrm{known}}(s_t, a)$ when probing frontiers; thus exploration focuses on unknown transitions.

\end{document}
